% ============================================================
% HW5 - Object Detection in Self-Driving Cars
% Fahad Ali Jalil | Universitat Klagenfurt | Summer 2026
% ============================================================
% To embed videos, compile with pdflatex.
% Required packages: media9 (for .avi embedding) or animate (for GIF)
% If media9 is unavailable, the fallback contact sheet images are shown.
% Compile command: pdflatex -shell-escape HW5_fahadjalil_kitti_yolo26.tex
% ============================================================

\documentclass[12pt, a4paper]{article}

% ── Packages ────────────────────────────────────────────────────────────────
\usepackage[a4paper, top=2.5cm, bottom=2.5cm, left=2.5cm, right=2.5cm]{geometry}
\usepackage{graphicx}
\usepackage{booktabs}
\usepackage{array}
\usepackage{tabularx}
\usepackage{longtable}
\usepackage{xcolor}
\usepackage{fancyhdr}
\usepackage{titlesec}
\usepackage{parskip}
\usepackage{microtype}
\usepackage{hyperref}
\usepackage{caption}
\usepackage{float}
\usepackage{tcolorbox}
\usepackage{fontenc}
\usepackage{inputenc}
\usepackage{lmodern}
\usepackage{setspace}
\usepackage{ragged2e}

% Video embedding — uses media9 for .avi, animate for GIFs
% Comment out if not installed and use fallback images instead
\usepackage{media9}
\usepackage{animate}

% ── Colours ──────────────────────────────────────────────────────────────────
\definecolor{darkblue}{HTML}{1a1a2e}
\definecolor{midblue}{HTML}{2c6b9e}
\definecolor{lightblue}{HTML}{ddeaf5}
\definecolor{annbg}{HTML}{eef2ff}
\definecolor{tablehead}{HTML}{1a1a2e}
\definecolor{tablerow1}{HTML}{f0f4ff}

% ── Section formatting ────────────────────────────────────────────────────────
\titleformat{\section}
  {\Large\bfseries\color{darkblue}}
  {\thesection.}{0.5em}{}
  [\color{midblue}\titlerule]

\titleformat{\subsection}
  {\normalsize\bfseries\color{midblue}}
  {\thesubsection}{0.5em}{}

\titlespacing{\section}{0pt}{14pt}{6pt}
\titlespacing{\subsection}{0pt}{10pt}{4pt}

% ── Annotation box ────────────────────────────────────────────────────────────
\tcbuselibrary{skins}
\newenvironment{annbox}{%
  \begin{tcolorbox}[
    colback=annbg,
    colframe=midblue,
    boxrule=0.5pt,
    arc=2pt,
    left=6pt, right=6pt, top=5pt, bottom=5pt,
    fontupper=\small\color{darkblue}
  ]
}{%
  \end{tcolorbox}
}

% ── Table styling ─────────────────────────────────────────────────────────────
\newcolumntype{L}[1]{>{\RaggedRight\arraybackslash}p{#1}}
\newcolumntype{C}[1]{>{\Centering\arraybackslash}p{#1}}

% Header row colour
\newcommand{\hdrrow}{\rowcolor{tablehead}\color{white}\bfseries}

% Alternating row colours
\usepackage{colortbl}
\newcommand{\rowA}{\rowcolor{tablerow1}}
\newcommand{\rowB}{\rowcolor{white}}

% ── Header and footer ─────────────────────────────────────────────────────────
\pagestyle{fancy}
\fancyhf{}
\renewcommand{\headrulewidth}{0pt}
\renewcommand{\footrulewidth}{0pt}

\fancyhead[L]{%
  \includegraphics[height=0.65cm]{images/logo_white.png}%
}
\fancyhead[R]{\small\color{gray} LAB on Autonomous Driving Cars \ | \ Fahad Ali Jalil}
\fancyfoot[L]{\small\color{gray} Universitat Klagenfurt \ | \ Summer 2026}
\fancyfoot[R]{\small\color{gray} Page \thepage}

% Blue rule under header
\renewcommand{\headrule}{{\color{midblue}\hrule width\headwidth height 1.2pt}}
% Blue rule above footer
\renewcommand{\footrule}{{\color{midblue}\hrule width\headwidth height 1.2pt}}

% ── Hyperlink setup ───────────────────────────────────────────────────────────
\hypersetup{
  colorlinks=true,
  linkcolor=midblue,
  urlcolor=midblue,
  pdftitle={HW5 - Object Detection in Self-Driving Cars},
  pdfauthor={Fahad Ali Jalil}
}

% ── Caption style ─────────────────────────────────────────────────────────────
\captionsetup{
  font=small,
  labelfont=it,
  textfont=it,
  justification=centering,
  skip=4pt
}

% ─────────────────────────────────────────────────────────────────────────────
\begin{document}

% ── COVER PAGE ───────────────────────────────────────────────────────────────
\thispagestyle{fancy}
\vspace*{2cm}

\begin{center}
  \includegraphics[width=7cm]{images/logo_white.png}\\[1.2cm]
  {\normalsize Homework 5}\\[0.3cm]
  {\LARGE\bfseries\color{darkblue}
    Object Detection in Self-Driving Cars\\
    using Python and Neural Networks}\\[0.5cm]
  {\color{midblue}\rule{\textwidth}{1.5pt}}
\end{center}

\vspace{0.6cm}

\begin{tabular}{@{}L{3.5cm} L{10cm}@{}}
  \textbf{Student}    & Fahad Ali Jalil \\[4pt]
  \textbf{Student ID} & 12502261 \\[4pt]
  \textbf{Course}     & LAB on Autonomous Driving Cars (700.385 / 700.386) \\[4pt]
  \textbf{Supervisor} & Dipl.-Ing. Mohamed Salem \\[4pt]
  \textbf{University} & Universitat Klagenfurt \\[4pt]
  \textbf{Dataset}    & KITTI 2D Object Detection \\[4pt]
  \textbf{Model}      & YOLO26n (Ultralytics, January 2026) \\[4pt]
  \textbf{Deadline}   & 15.06.2026 \\
\end{tabular}

\newpage

% ─────────────────────────────────────────────────────────────────────────────
\section{Task Description}

The objective of Homework 5 is to implement an end-to-end object detection pipeline for autonomous
driving using Python and deep learning. The task requires selecting a dataset, preprocessing it,
choosing and implementing a neural network architecture, training the model, evaluating performance
on a validation set, and visualising detection results on sample images.

\vspace{0.3cm}
\begin{tabular}{@{} C{1.2cm} L{12.5cm} @{}}
  \hdrrow Step & \hdrrow Requirement \\
  \rowA 1 & Choose a well-known dataset: CityScapes, KITTI, or similar \\
  \rowB 2 & Preprocess: resize images, normalise pixels, handle annotations, split train/val \\
  \rowA 3 & Choose an architecture: Faster R-CNN, YOLO, or SSD \\
  \rowB 4 & Implement using Python and a deep learning framework \\
  \rowA 5 & Train the model and evaluate on the validation set \\
  \rowB 6 & Visualise bounding box detections on sample images \\
\end{tabular}

% ─────────────────────────────────────────────────────────────────────────────
\section{Approach}

The Ultralytics YOLO26 framework was selected. YOLO26n (January 2026) is the latest YOLO
generation featuring a fully end-to-end NMS-free architecture. All six homework requirements are
satisfied.

\vspace{0.3cm}
\begin{tabular}{@{} L{4.5cm} L{9.2cm} @{}}
  \hdrrow Homework Step & \hdrrow How it is fulfilled \\
  \rowA Dataset           & KITTI downloaded automatically via kitti.yaml (390.5 MB) \\
  \rowB Resize images     & All images resized to 640$\times$640 pixels during data loading \\
  \rowA Normalise pixels  & Pixel values scaled from 0--255 to 0.0--1.0 by the dataloader \\
  \rowB Annotations       & kitti.yaml maps each label file to its corresponding image \\
  \rowA Train/val split   & 5,985 training images and 1,496 validation images \\
  \rowB Architecture      & YOLO26n: 260 layers, 2,506,920 parameters, 5.8 GFLOPs \\
  \rowA Loss function     & box\_loss + cls\_loss + dfl\_loss (computed each epoch) \\
  \rowB Training pipeline & model.train() with AdamW optimizer, 10 epochs, batch size 16 \\
  \rowA Evaluation        & mAP50 and mAP50-95 computed on 1,496 validation images \\
  \rowB Visualisation     & model.predict() with save=True draws boxes on sample images \\
\end{tabular}

% ─────────────────────────────────────────────────────────────────────────────
\section{Dataset: KITTI}

The KITTI Vision Benchmark Suite contains real-world driving scenes collected in Karlsruhe, Germany.
It contains 8 annotated object classes covering all road users encountered in urban and highway driving.

\vspace{0.3cm}
\begin{tabular}{@{} L{4.5cm} L{9.2cm} @{}}
  \hdrrow Property & \hdrrow Value \\
  \rowA Total images      & 7,481 \\
  \rowB Training images   & 5,985 \\
  \rowA Validation images & 1,496 \\
  \rowB Input resolution  & Variable original, resized to 640$\times$640 for training \\
  \rowA Download size     & 390.5 MB \\
  \rowB Classes           & car, van, truck, pedestrian, person\_sitting, cyclist, tram, misc \\
  \rowA Source            & Geiger et al., IJRR 2013 \\
\end{tabular}

\vspace{0.4cm}
\subsection{Figure 1: Class Distribution (Training Set)}

\begin{figure}[H]
  \centering
  \includegraphics[width=0.65\textwidth]{images/label_bar_only.png}
  \caption{Annotation count per class in the KITTI training set.}
\end{figure}

\begin{annbox}
Car dominates with 23,026 annotations out of $\sim$32,437 total (71\%). Rare classes such as
person\_sitting (186) and tram (394) are underrepresented, explaining their lower detection accuracy
in Section 6. Box centres cluster at y=0.5--0.6 (lower half of frame), consistent with road-level
objects in a forward-facing camera.
\end{annbox}

\newpage

% ─────────────────────────────────────────────────────────────────────────────
\section{Model: YOLO26n}

YOLO26n is the nano (smallest) variant of YOLO26, pre-trained on COCO (80 classes) and fine-tuned
on KITTI (8 classes). It outputs up to 300 detections per image without requiring NMS post-processing.

\vspace{0.3cm}
\begin{tabular}{@{} L{4.5cm} L{9.2cm} @{}}
  \hdrrow Property & \hdrrow Value \\
  \rowA Layers          & 260 \\
  \rowB Parameters      & 2,506,920 \\
  \rowA GFLOPs          & 5.8 \\
  \rowB Input size      & 640$\times$640 pixels \\
  \rowA Output          & (1, 300, 6) --- NMS-free, max 300 detections per image \\
  \rowB Pretrained on   & COCO dataset (80 classes) \\
  \rowA Fine-tuned on   & KITTI dataset (8 classes) \\
  \rowB Optimizer       & AdamW (lr=0.000833, momentum=0.9, weight decay=0.0005) \\
  \rowA GPU             & Tesla T4, 14,913 MiB --- Google Colab \\
\end{tabular}

\vspace{0.3cm}
\begin{tabular}{@{} L{2.5cm} L{4.0cm} L{7.2cm} @{}}
  \hdrrow Loss & \hdrrow Full Name & \hdrrow What it penalises \\
  \rowA box\_loss & Bounding Box Loss       & Distance between predicted and ground-truth box coordinates \\
  \rowB cls\_loss & Classification Loss     & Error in predicting the correct class label \\
  \rowA dfl\_loss & Distribution Focal Loss & Precision of predicted box edge positions \\
\end{tabular}

% ─────────────────────────────────────────────────────────────────────────────
\section{Training Progress}

The model was trained for 10 epochs on a Tesla T4 GPU. All loss values decrease consistently and
mAP50 improves from approximately 0.19 at epoch 1 to 0.653 at epoch 10.

\vspace{0.3cm}
\begin{tabular}{@{} C{1.5cm} C{1.8cm} C{1.8cm} C{1.8cm} C{1.8cm} C{1.8cm} C{2cm} @{}}
  \hdrrow Epoch & \hdrrow GPU Mem & \hdrrow box\_loss & \hdrrow cls\_loss & \hdrrow dfl\_loss & \hdrrow mAP50 & \hdrrow mAP50-95 \\
  \rowA 1/10  & 2.44G & 1.767 & 4.703 & 0.006 & $\sim$0.19 & $\sim$0.14 \\
  \rowB 8/10  & 2.34G & 1.083 & 0.788 & 0.998 & 0.613 & 0.385 \\
  \rowA 9/10  & 2.36G & 1.044 & 0.744 & 0.985 & 0.622 & 0.391 \\
  \rowB 10/10 & 2.36G & 1.010 & 0.710 & 0.971 & 0.653 & 0.421 \\
\end{tabular}

\vspace{0.4cm}
\subsection{Figure 2: Training and Validation Losses and Metrics}

\begin{figure}[H]
  \centering
  \includegraphics[width=\textwidth]{images/training.png}
  \caption{Top row: training losses and validation precision/recall across 10 epochs.
           Bottom row: validation losses and mAP scores.}
\end{figure}

\begin{annbox}
box\_loss drops from 1.77 to 1.30 --- localisation improves steadily. cls\_loss drops sharply from 4.0
to 1.0 in the first 4 epochs --- class recognition learned quickly. mAP50 and mAP50-95 trend upward
with no sign of overfitting within 10 epochs, indicating further improvement with additional training.
\end{annbox}

\newpage

% ─────────────────────────────────────────────────────────────────────────────
\section{Evaluation Results}

The best checkpoint was evaluated on 1,496 validation images across all 8 KITTI classes.

\vspace{0.3cm}
\begin{tabular}{@{} L{3.0cm} C{1.5cm} C{1.8cm} C{2.0cm} C{1.6cm} C{1.6cm} C{1.8cm} @{}}
  \hdrrow Class & \hdrrow Images & \hdrrow Instances & \hdrrow Precision & \hdrrow Recall & \hdrrow mAP50 & \hdrrow mAP50-95 \\
  \rowA ALL (overall)  & 1496 & 8128 & 0.727 & 0.571 & 0.653 & 0.422 \\
  \rowB car            & 1322 & 5716 & 0.882 & 0.818 & 0.902 & 0.661 \\
  \rowA van            & 419  & 555  & 0.819 & 0.589 & 0.720 & 0.517 \\
  \rowB truck          & 212  & 223  & 0.939 & 0.763 & 0.852 & 0.656 \\
  \rowA pedestrian     & 377  & 921  & 0.693 & 0.481 & 0.578 & 0.283 \\
  \rowB person\_sitting & 17  & 36   & 0.425 & 0.361 & 0.414 & 0.206 \\
  \rowA cyclist        & 252  & 352  & 0.633 & 0.484 & 0.526 & 0.269 \\
  \rowB tram           & 77   & 117  & 0.684 & 0.658 & 0.757 & 0.474 \\
  \rowA misc           & 158  & 208  & 0.741 & 0.412 & 0.475 & 0.310 \\
\end{tabular}

\vspace{0.4cm}
\subsection{Figure 3: Confusion Matrix (Raw Counts)}

\begin{figure}[H]
  \centering
  \includegraphics[width=0.85\textwidth]{images/cm.png}
  \caption{Raw confusion matrix. Rows = predicted class, columns = true class. Background row = missed objects.}
\end{figure}

\begin{annbox}
Car: 4,362 correct detections out of 5,716 --- strongest performer. 1,326 cars missed (background row),
mainly small or distant vehicles. Pedestrian: 548 missed --- small objects at driving range are harder to detect.
Person\_sitting and cyclist frequently missed due to low training count and small visual size.
\end{annbox}

\vspace{0.3cm}
\subsection{Figure 4: Confusion Matrix (Normalised)}

\begin{figure}[H]
  \centering
  \includegraphics[width=0.85\textwidth]{images/cm_norm.png}
  \caption{Normalised confusion matrix --- fraction of ground-truth instances predicted as each class.}
\end{figure}

\begin{annbox}
Car: 0.76 true positive rate. Truck: 0.68. Van: 0.50 --- frequently confused with car (0.09) due to
visual similarity. Pedestrian: 0.39 --- 60\% missed as background at driving distances.
Person\_sitting: 0.31 --- hardest class with only 36 validation instances.
\end{annbox}

\newpage

% ─────────────────────────────────────────────────────────────────────────────
\section{Bounding Box Visualisation on Sample Images}

The trained model was applied to 6 images: one fixed test image and 5 randomly selected KITTI
validation images (random seed 42). All detection images are 1242$\times$375 pixels.

\subsection{Sample 1: Urban Street --- Fixed Test Image}
\begin{figure}[H]
  \centering
  \includegraphics[width=\textwidth]{images/det1.png}
  \caption{Sample 1: Dense urban street --- 12 cars and 1 truck detected.}
\end{figure}
\begin{annbox}
12 cars detected (confidence 0.26--0.97) and 1 truck (0.65). Large nearby cars on both edges detected
at high confidence (0.89, 0.88). Multiple overlapping cars in the centre are individually separated.
Partially visible car on far left detected at 0.26 --- expected for a heavily occluded object.
\end{annbox}

\subsection{Sample 2: Quiet Residential Street}
\begin{figure}[H]
  \centering
  \includegraphics[width=\textwidth]{images/det2.png}
  \caption{Sample 2: Quiet residential street --- 1 car detected at a distance (confidence 0.88).}
\end{figure}
\begin{annbox}
1 car detected at 0.88 confidence at significant distance down the road. Scene contains many
distractors --- hedges, fences, lamp posts --- yet the model identifies the single vehicle correctly
and produces zero false positives.
\end{annbox}

\subsection{Sample 3: Residential Parking Area}
\begin{figure}[H]
  \centering
  \includegraphics[width=\textwidth]{images/det3.png}
  \caption{Sample 3: Residential parking area --- 5 cars detected.}
\end{figure}
\begin{annbox}
5 cars: 0.46 (left edge, partially cut off), 0.71 (white parked car), 0.95 (blue car in foreground),
0.95 (dark car right-centre), 0.95 (partially visible far right). Truncated frame-edge car scores
lower (0.46) as expected.
\end{annbox}

\newpage

\subsection{Sample 4: Tree-Lined Suburban Road}
\begin{figure}[H]
  \centering
  \includegraphics[width=\textwidth]{images/det4.png}
  \caption{Sample 4: Tree-lined suburban road --- 2 cars detected at high confidence.}
\end{figure}
\begin{annbox}
Silver estate car on the left (0.88) and dark car ahead on the road (0.94). Distant dark car is small
in frame but detected at 0.94 --- confirming correct detection at driving distances. Tram tracks on
the right produce no false detections.
\end{annbox}

\subsection{Sample 5: Rural Highway with Oncoming Traffic}
\begin{figure}[H]
  \centering
  \includegraphics[width=\textwidth]{images/det5.png}
  \caption{Sample 5: Rural highway --- 4 cars detected, tram warning sign correctly ignored.}
\end{figure}
\begin{annbox}
4 cars: partially visible dark car at left edge (0.91), red car beside railway (0.89), blue car ahead
(0.95), black car immediately ahead (0.92). All scores above 0.89. Tram-crossing warning sign is
correctly not detected --- road signs are not a KITTI class.
\end{annbox}

\subsection{Sample 6: Narrow Cobblestone Street}
\begin{figure}[H]
  \centering
  \includegraphics[width=\textwidth]{images/det6.png}
  \caption{Sample 6: Narrow cobblestone street --- 4 cars detected in tight cluster.}
\end{figure}
\begin{annbox}
4 cars in right half: 0.95, 0.82, 0.51, 0.95. Fifth car at far right edge at 0.76. Three tightly
clustered cars around garage entrance individually distinguished. Large empty left half of frame
(metal fence, building wall) produces zero false positives.
\end{annbox}

\newpage

% ─────────────────────────────────────────────────────────────────────────────
\section{Model Export}

The trained model was exported to ONNX format for deployment on any hardware without requiring PyTorch.

\vspace{0.3cm}
\begin{tabular}{@{} L{4.5cm} L{9.2cm} @{}}
  \hdrrow Property & \hdrrow Value \\
  \rowA Format       & ONNX (Open Neural Network Exchange) \\
  \rowB File size    & 9.4 MB \\
  \rowA ONNX opset   & 17 (re-exported for Python 3.7 compatibility) \\
  \rowB Input shape  & (1, 3, 640, 640) \\
  \rowA Output shape & (1, 300, 6) --- up to 300 detections per image \\
  \rowB Saved to     & /home/fahad/carla\_0.9.11/kitti\_yolo26n/best.onnx \\
\end{tabular}

% ─────────────────────────────────────────────────────────────────────────────
\section{Model Training Summary}

After 10 epochs the best checkpoint achieved overall mAP50 of 0.653 on the KITTI validation set.
Performance is strongest for large, frequent classes --- car (0.902), truck (0.852), van (0.720) ---
and lower for rare, small classes --- person\_sitting (0.414), cyclist (0.526). Inference runs at
approximately 13 ms per image on a T4 GPU, suitable for real-time deployment.

\newpage

% ─────────────────────────────────────────────────────────────────────────────
\section{CARLA Simulation with Live YOLO26n Detection}

After successfully training and evaluating the YOLO26n model on the KITTI dataset, the trained model
was tested inside CARLA (Car Learning to Act) --- an open-source autonomous driving simulator. CARLA
generates realistic camera images from a virtual vehicle, making it a natural environment to verify
whether a model trained on real photos can also detect objects in a simulated world.

\subsection{Approach and Thought Process}

The idea is straightforward: mount a virtual camera on a manually driven vehicle in CARLA, feed every
camera frame into the trained YOLO26n model, and draw bounding boxes around detected objects in real
time. This tests \textit{domain transfer} --- whether a model trained on real German driving photos
(KITTI) can detect objects in a simulated world.

A separate Python script was created for this task: \texttt{carla\_yolo26\_detection.py}. This script
connects to the CARLA server, spawns the player vehicle and NPC traffic, attaches a camera sensor,
runs YOLO26n inference on each frame, draws detections on the live feed, and saves video recordings
on demand using the R key.

A key technical challenge was that the trained model requires Python 3.8+ while CARLA 0.9.11 only
provides a Python 3.7 API. The solution was to export the model to ONNX format --- a universal format
loadable by any Python version via \texttt{onnxruntime} --- allowing the full pipeline to run in
Python 3.7 alongside CARLA with no version conflicts.

\subsection{Implementation Steps}

\vspace{0.2cm}
\begin{tabular}{@{} C{0.8cm} L{13cm} @{}}
  \hdrrow Step & \hdrrow What was done \\
  \rowA 1 & Trained model exported to ONNX format (opset 17) for Python 3.7 compatibility \\
  \rowB 2 & Conda environment carla37 set up with Python 3.7, CARLA egg, onnxruntime, opencv \\
  \rowA 3 & carla\_yolo26\_detection.py written with CARLA connection, camera, NPC, detection, recording \\
  \rowB 4 & CARLA launched with \texttt{-quality-level=Low} to reduce GPU load on integrated Intel Iris Xe GPU \\
  \rowA 5 & Town03 loaded with 15 NPC vehicles and 6 pedestrian walkers \\
  \rowB 6 & Vehicle driven manually with W/A/S/D, live detections shown on camera window \\
  \rowA 7 & R key used to record sessions --- saved as timestamped .avi files \\
\end{tabular}

\subsection{Figure 8: CARLA Launch --- Low Quality Mode}

\begin{figure}[H]
  \centering
  \includegraphics[width=\textwidth]{images/carla_low.png}
  \caption{CARLA 0.9.11 launched with \texttt{-quality-level=Low} flag to reduce GPU and memory usage.}
\end{figure}
\begin{annbox}
CARLA launched using \texttt{-quality-level=Low} which reduces texture resolution, shadow quality, and
post-processing. Necessary because the test machine uses an integrated Intel Iris Xe GPU with shared
memory rather than a dedicated GPU.
\end{annbox}

\subsection{Figure 9: Successful Connection and Environment Load}

\begin{figure}[H]
  \centering
  \includegraphics[width=0.9\textwidth]{images/carla_env.png}
  \caption{Terminal output confirming CARLA connected, Town03 loaded, vehicle and camera spawned.}
\end{figure}
\begin{annbox}
The script connected to CARLA at localhost:2000, loaded Town03, spawned the player vehicle (Tesla
Model 3), attached the RGB camera at 800$\times$600, and spawned 15 NPC vehicles and 6 walkers. Font
warnings in the terminal are harmless Qt messages that do not affect the simulation.
\end{annbox}

\subsection{Figure 10: Live Detection Window Running in CARLA}

\begin{figure}[H]
  \centering
  \includegraphics[width=\textwidth]{images/carla_running.png}
  \caption{Live YOLO26n detection window showing CARLA camera feed with bounding boxes.}
\end{figure}
\begin{annbox}
The detection window shows the forward-facing camera view from the player vehicle in Town03. Speed and
gear indicator appear top-left. Controls shown at the bottom. The model runs ONNX inference on every
frame and draws colour-coded boxes: green for cars, red for pedestrians, yellow for cyclists. The REC
indicator appears top-right when recording is active.
\end{annbox}

\newpage

\subsection{Domain Gap Observation}

The model was trained on real KITTI photographs from Karlsruhe, Germany. CARLA generates synthetic
images rendered by a game engine --- this visual difference is called the \textit{domain gap}. In
practice the model still detects cars and pedestrians in CARLA, but with slightly lower confidence
scores than on real KITTI images. This is a known and expected challenge in autonomous driving research.

\subsection{Recorded Driving Sessions}

Two driving sessions were recorded using the R key toggle. Three representative frames from each
session are shown below. The full video files can be played directly in this PDF using the embedded
player below each contact sheet (requires Adobe Acrobat or a compatible PDF reader).

% ── RECORDING 1 ──────────────────────────────────────────────────────────────
\vspace{0.3cm}
\subsubsection*{Recording 1 --- carla\_yolo26n\_20260614\_152239.avi (3.5 seconds)}

\begin{figure}[H]
  \centering
  \includegraphics[width=\textwidth]{images/gif1_sheet.png}
  \caption{Recording 1: Three frames sampled from the first driving session (beginning, middle, end).}
\end{figure}

\begin{annbox}
First recording session --- 3.5 second clip, 35 frames at 20 FPS, resolution 800$\times$600. Shows
the initial moments of the driving session in Town03. YOLO26n detects road objects frame by frame with
colour-coded bounding boxes.
\end{annbox}

\vspace{0.3cm}
% Embedded video --- requires Adobe Acrobat to play
% If media9 is not available, comment out the includemedia block below
\begin{center}
\includemedia[
  width=0.85\textwidth,
  height=0.4\textwidth,
  activate=pageopen,
  passcontext,
  addresource=images/carla_yolo26n_20260614_152239.avi,
  flashvars={
    source=images/carla_yolo26n_20260614_152239.avi
    &autoPlay=false
    &loop=true
  }
]{%
  \includegraphics[width=0.85\textwidth]{images/gif1_sheet.png}
}{VPlayer.swf}
\end{center}

\vspace{0.2cm}
\begin{center}
  \small\textit{$\blacktriangleright$ Click to play Recording 1 (requires Adobe Acrobat Reader)}
\end{center}

\newpage

% ── RECORDING 2 ──────────────────────────────────────────────────────────────
\subsubsection*{Recording 2 --- carla\_yolo26n\_20260614\_152530.avi (27 seconds)}

\begin{figure}[H]
  \centering
  \includegraphics[width=\textwidth]{images/gif2_sheet.png}
  \caption{Recording 2: Three frames sampled from the second driving session (beginning, middle, end).}
\end{figure}

\begin{annbox}
Second recording session --- 27 second clip, 539 frames at 20 FPS, resolution 800$\times$600. Longer
driving sequence through Town03 with NPC vehicles and pedestrians present. Green boxes show detected
cars, red boxes show detected pedestrians. Confidence scores are shown above each bounding box.
\end{annbox}

\vspace{0.3cm}
% Embedded video --- requires Adobe Acrobat to play
\begin{center}
\includemedia[
  width=0.85\textwidth,
  height=0.4\textwidth,
  activate=pageopen,
  passcontext,
  addresource=images/carla_yolo26n_20260614_152530.avi,
  flashvars={
    source=images/carla_yolo26n_20260614_152530.avi
    &autoPlay=false
    &loop=true
  }
]{%
  \includegraphics[width=0.85\textwidth]{images/gif2_sheet.png}
}{VPlayer.swf}
\end{center}

\vspace{0.2cm}
\begin{center}
  \small\textit{$\blacktriangleright$ Click to play Recording 2 (requires Adobe Acrobat Reader)}
\end{center}

\newpage

% ─────────────────────────────────────────────────────────────────────────────
\section{Conclusion}

This homework implemented a complete object detection pipeline for autonomous driving using YOLO26n
fine-tuned on the KITTI dataset. All six homework requirements were fulfilled. After 10 epochs on a
Google Colab T4 GPU, the model achieved an overall mAP50 of 0.653, with car reaching 0.902, truck
0.852, and van 0.720. The model was then successfully deployed inside the CARLA simulator where it
detected vehicles and pedestrians in real time on a manually driven vehicle --- demonstrating practical
applicability beyond static dataset evaluation. A separate Python script handles the full CARLA
integration pipeline including ONNX inference, live HUD overlay, and video recording. Two driving
sessions were recorded as documentation.

% ─────────────────────────────────────────────────────────────────────────────
\section*{References}
\addcontentsline{toc}{section}{References}

\begin{enumerate}
  \item Geiger, A., Lenz, P., Stiller, C., and Urtasun, R. (2013). Vision meets Robotics: The KITTI
        Dataset. \textit{International Journal of Robotics Research (IJRR)}.
  \item Jocher, G., and Qiu, J. (2026). Ultralytics YOLO26: Unified Real-Time End-to-End Vision
        Models. \textit{arXiv:2606.03748}.
  \item Dosovitskiy, A. et al. (2017). CARLA: An Open Urban Driving Simulator. \textit{CoRL 2017}.
  \item Ultralytics YOLO26 Documentation.
        \url{https://docs.ultralytics.com/models/yolo26/}
  \item KITTI Dataset on Ultralytics Docs.
        \url{https://docs.ultralytics.com/datasets/detect/kitti/}
\end{enumerate}

\end{document}
